\section{振幅増幅}
\label{sec:amplitude-amplification}

\subsection{問題の設定}
\label{subsec:amplitude-amplification-setting}

量子アルゴリズム$\mathcal{A}$が$n$量子ビットの初期状態から
\begin{equation}
  \ket{\psi}
  \coloneq
  \mathcal{A}\ket{0}^{\otimes n}
  \label{eq:amplitude-amplification-initial-state}
\end{equation}
を生成するとする．
計算基底状態の集合を，目的の条件を満たす\textbf{良い状態}と，
条件を満たさない\textbf{悪い状態}に分ける．
良い状態が張る部分空間への射影演算子を$\Pi_{\mathrm{g}}$とし，
初期状態を測定したときに良い状態を得る確率を
\begin{equation}
  a
  \coloneq
  \bra{\psi}\Pi_{\mathrm{g}}\ket{\psi}
  \label{eq:initial-good-probability}
\end{equation}
と定義する．

成功確率が$0<a<1$を満たすとし，規格化した良い成分と悪い成分をそれぞれ
\begin{align}
  \ket{\psi_{\mathrm{g}}}
  &\coloneq
  \frac{\Pi_{\mathrm{g}}\ket{\psi}}{\sqrt{a}},
  \label{eq:normalized-good-state}\\
  \ket{\psi_{\mathrm{b}}}
  &\coloneq
  \frac{(I-\Pi_{\mathrm{g}})\ket{\psi}}{\sqrt{1-a}}
  \label{eq:normalized-bad-state}
\end{align}
と定義する．
さらに
\begin{equation}
  \theta
  \coloneq
  \arcsin\sqrt{a},
  \qquad
  0<\theta<\frac{\pi}{2}
  \label{eq:amplification-angle}
\end{equation}
とおくと，初期状態は
\begin{equation}
  \ket{\psi}
  =
  \cos\theta\ket{\psi_{\mathrm{b}}}
  +\sin\theta\ket{\psi_{\mathrm{g}}}
  \label{eq:good-bad-decomposition}
\end{equation}
と表される．

同じ量子回路を実行して直ちに測定する方法では，
良い状態を得るまでに平均して$O(1/a)$回の試行が必要になる．
\textbf{振幅増幅}は，測定前に良い成分の確率振幅を増幅し，
必要な反復回数を$O(1/\sqrt{a})$へ減らす方法である\cite{shimada2020}．

\subsection{良い部分空間に関する反射}
\label{subsec:good-subspace-reflection}

良い状態と悪い状態を判定する量子オラクルを用いて，
良い部分空間の位相だけを反転する演算子を
\begin{equation}
  S_{\mathrm{g}}
  \coloneq
  I-2\Pi_{\mathrm{g}}
  \label{eq:good-state-reflection}
\end{equation}
と定義する．
この演算子は
\begin{align*}
  S_{\mathrm{g}}\ket{\psi_{\mathrm{b}}}
  &=
  \ket{\psi_{\mathrm{b}}},\\
  S_{\mathrm{g}}\ket{\psi_{\mathrm{g}}}
  &=
  -\ket{\psi_{\mathrm{g}}}
\end{align*}
と作用する．
したがって，$S_{\mathrm{g}}$は良い部分空間へ符号反転の目印を付ける操作であり，
幾何学的には悪い部分空間を軸とする反射である．

判定関数$\chi(x)\in\{0,1\}$が利用できる場合，位相オラクルを
\begin{equation}
  O_{\chi}\ket{x}
  \coloneq
  (-1)^{\chi(x)}\ket{x}
  \label{eq:phase-oracle}
\end{equation}
として構成できる．
判定値$\chi(x)=1$を良い状態とすれば，$O_{\chi}$は$S_{\mathrm{g}}$に対応する．

\subsection{初期状態に関する反射}
\label{subsec:initial-state-reflection}

初期状態$\ket{\psi}$を軸とする反射演算子を
\begin{equation}
  S_{\psi}
  \coloneq
  2\ket{\psi}\bra{\psi}-I
  \label{eq:initial-state-reflection}
\end{equation}
と定義する．
状態準備の関係$\ket{\psi}=\mathcal{A}\ket{0}^{\otimes n}$より，
\begin{equation}
  S_{\psi}
  =
  \mathcal{A}
  \left(
    2\ket{0}^{\otimes n}\bra{0}^{\otimes n}-I
  \right)
  \mathcal{A}^{\dagger}
  \label{eq:prepared-state-reflection}
\end{equation}
と表される．
したがって，状態準備回路$\mathcal{A}$，その逆回路$\mathcal{A}^{\dagger}$，
および$\ket{0}^{\otimes n}$だけを位相反転する回路を用いて$S_{\psi}$を構成できる．

\subsection{二つの反射による回転}
\label{subsec:two-reflections-as-rotation}

一回の振幅増幅反復を表す演算子を
\begin{equation}
  Q
  \coloneq
  S_{\psi}S_{\mathrm{g}}
  \label{eq:amplitude-amplification-operator}
\end{equation}
と定義する．
順序付き基底$\{\ket{\psi_{\mathrm{b}}},\ket{\psi_{\mathrm{g}}}\}$に関して表すと
\begin{equation}
  Q
  =
  \begin{pmatrix}
    \cos 2\theta & -\sin 2\theta\\
    \sin 2\theta & \cos 2\theta
  \end{pmatrix}
  \label{eq:amplitude-amplification-rotation-matrix}
\end{equation}
となる．
すなわち，二つの反射を続けて作用させると，
良い状態と悪い状態が張る二次元平面上で角度$2\theta$の回転が生じる\cite{shimada2020}．

初期状態は悪い状態の軸から角度$\theta$だけ良い状態側へ傾いている．
したがって，$Q$を$r$回作用させると
\begin{equation}
  Q^r\ket{\psi}
  =
  \cos((2r+1)\theta)\ket{\psi_{\mathrm{b}}}
  +\sin((2r+1)\theta)\ket{\psi_{\mathrm{g}}}
  \label{eq:amplitude-after-r-iterations}
\end{equation}
となる．
この状態を測定して良い結果を得る確率は
\begin{equation}
  p_r
  \coloneq
  \sin^2((2r+1)\theta)
  \label{eq:amplified-success-probability}
\end{equation}
である．
一回の反復で角度が$2\theta$ずつ増加し，状態が良い部分空間へ近づくことが分かる．

\subsection{反復回数と二次高速化}
\label{subsec:amplification-iteration-count}

成功確率を最大にするには，$(2r+1)\theta$が$\pi/2$に最も近くなるように$r$を選ぶ．
したがって，必要な反復回数はおよそ
\begin{equation}
  r
  \simeq
  \frac{\pi}{4\theta}-\frac{1}{2}
  \label{eq:near-optimal-amplification-count}
\end{equation}
である．
成功確率$a$が小さい場合には$\theta=\arcsin\sqrt{a}\simeq\sqrt{a}$であるため，
反復回数は$O(1/\sqrt{a})$となる．

一方，最適な回数を超えて$Q$を作用させると，状態は良い部分空間を通り過ぎ，
成功確率は再び低下する．
したがって，初期成功確率$a$が未知の場合には，反復回数を固定するだけでは安定して増幅できない．
本論文では$a$が既知である理想的な場合の反復回数を用い，
未知の場合の反復方法には立ち入らない．
